\problem{Aimantation d'une substance paramagnétique}

\question{Montrez que l'on peut obtenir $\vb{M} = \overline{\bm{\mu}_{\mathrm{tot}}}/{V}$ à partir de la formule générale $\vb{M} = \frac{1}{\beta V} \pdv{\ln Z}{\vb{B}}$.}
Commençons de la formule générale du moment dipolaire magnétique moyen par unité de volume:
\[
    \vb{M} &= \frac{1}{\beta V} \pdv{\ln Z}{\vb{B}} \notag \\
    &= \frac{1}{\beta V} \frac{1}{Z} \pdv{Z}{\vb{B}} \notag
    \intertext{Or, on sait que $Z = \sum_{r} e^{-\beta\epsilon_r}$, d'où:} 
    \vb{M} &= \frac{1}{\beta V} \frac{1}{Z} \pdv{\vb{B}} \qty(\sum_r e^{\beta\mu_r\vb{B}})\notag \\
    &= \frac{1}{V Z} \sum_r \mu_r \pdv{\vb{B}} \qty(e^{\beta\mu_r\vb{B}}) \notag \\
    &= \frac{1}{V} \sum_r \mu_r \frac{e^{-\beta\epsilon_r}}{Z} \notag \\
    \intertext{On sait que $P_r = \frac{e^{-\beta\epsilon_r}}{Z}$. On remplaçe pour trouver:}
    &= \frac{1}{V} \sum_r \mu_r P_r. \notag \\
    \intertext{Sachant que la valeur moyenne d'une quantité $O$ est $\overline{O} = \sum_{i} O_i P_i$, on peut trouver:}
    \vb{M} &= \frac{\overline{\bm{\mu}}}{V}.
\]

\question{Montrez que $\overline{E} = -V\vb{M}\vb{M}$.}
Par définiton, on a l'énergie moyenntel que:
\[
    \overline{\epsilon} &= \sum_r \epsilon_r P_r. \notag
\]
Sachant que 
\[
    \implies \overline{\epsilon} &= -\sum_r \vb{B}\mu_r P_r \notag \\
    \iff &= - \vb{B}\overline{\mu} + \dots \notag \\
    &= -\vb{B}\vb{M}V + \dots, \label{eq: 2.1}
\]
puisque $\overline{\mu} = V \vb{M}$.

\question{Z pour N particules de spin 1?}
On sait que $Z_N = Z_1^N$. Trouvons $Z_1$:
\[
    Z_1 &= \sum_r e^{-\beta\epsilon_r} \notag \\
    &= \sum_r e^{-\beta\mu_B\vb{B}} \notag \\
    &= \sum_r \qty[e^{\beta\mu_B\vb{B}} + 1 + e^{-\beta\mu_B\vb{B}}] \notag \\
    &= 1 + \cosh\qty(\beta\mu_B\vb{B}). \notag
\]
Trouvons maintenant $Z_N$:
\[
    Z_N = \qty[1 + 2\cosh\qty(\beta\mu_B\vb{B})]^N. \label{eq: 2.2}
\]

\question{Calculez l'aimantation du système.}
En substituant l'équation \eqref{eq: 2.2} dans \eqref{eq: 2.1} on obtient:
\[
    \vb{M} &= \frac{1}{\beta V} \pdv{\ln Z_1^N}{\vb{B}}\notag \\
    &= \frac{N}{Z_1 \beta V } pdv{Z_1}{\vb{B}} \notag \\
    &= \frac{N}{Z_1 \beta V } \pdv{\vb{B}}\qty[e^{\beta\mu_B\vb{B}} + 1 + e^{-\beta\mu_B\vb{B}}] \notag \\
    &= \frac{\mu_B N}{V} \frac{2\sinh\qty(\beta\mu_B\vb{B})}{1 + 2\cosh\qty(\beta\mu_B\vb{B})}.
\]

\question{}
La température caractéristique est définie telle que:
\[
    k_B T &\approx \mu_B \vb{B}\notag \\
    \implies T_c &\approx \frac{\mu_B\vb{B}}{k_B}
\]
avec
\[
    \mu_B &= 9.27 * 10_{-24} \mathrm{J/T} \notag \\
    k_B &= 1.38 * 10_{-23} \mathrm{J/K}.
\]

on peux alors trouver $T_c$ tel que:
\[
    T_c \approx 0.67 \vb{B} \mathrm{J/T}.
\]

La susceptibilité magnétique linéaire est définie telle que:
\[
    \chi_M = \frac{\vb{M}}{\vb{B}}. 
\]

à haute température, on a:
\[
\beta\mu_B\vb{B} &\gg 1\notag \\
\implies &\begin{cases}
    \sinh(\beta\mu_B\vb{B}) \approx \beta\mu_B\vb{B}\\
    \sinh(\beta\mu_B\vb{B}) \approx 1
\end{cases} \notag \\
\implies \vb{M} &\approx \frac{2\beta}{3}\frac{\mu_B^2N\vb{B}}{V} \notag \\
\implies \chi_M &\approx \frac{2\beta}{3}\frac{\mu_B^2N}{V}.
\]
on trouve que 
\[
    \chi \propto \frac{1}{T}.
\]

À basse trempérature, on a:
\[
\beta\mu_B\vb{B} &\ll 1\notag \\
\implies &\begin{cases}
    \sinh(\beta\mu_B\vb{B}) \approx \frac{e^{\beta\mu_B\vb{B}}}{2}\\
    \sinh(\beta\mu_B\vb{B}) \approx \frac{e^{\beta\mu_B\vb{B}}}{2}
\end{cases} \notag \\
\implies \vb{M} &\approx \frac{N\mu_B}{V} \notag \\
\implies \chi_M &\approx \frac{N\mu_B}{V}.
\]

Ici, la magnétisation atteint un maximum. donc, tous les spins sont alignés avec le champ $\vb{B}$.

\clearpage
